% 55e Olympiade Internationale de Mathématiques — Le Cap, Afrique du Sud.
% Énoncés officiels en français (imo-official.org). Chaque problème vaut 7 points.

\begin{center}
  \emph{Premier jour --- mardi 8 juillet 2014 --- Durée : 4 heures 30}
\end{center}

\subsection*{Problème 1}
Soit $a_0 < a_1 < a_2 < \cdots$ une suite infinie d'entiers strictement positifs. Prouver qu'il existe un
unique entier $n \geq 1$ tel que
\[
a_n < \frac{a_0 + a_1 + \cdots + a_n}{n} \leq a_{n+1}.
\]

\subsection*{Problème 2}
Soit $n \geq 2$ un entier. On considère un échiquier $n \times n$ divisé en $n^2$ cases. Une configuration de
$n$ jetons répartis dans ces cases est dite \emph{paisible} si chaque ligne et chaque colonne de l'échiquier
contient exactement un jeton. Déterminer le plus grand entier strictement positif $k$ tel que, pour toute
configuration paisible de $n$ jetons, il existe un carré $k \times k$ qui ne contient aucun jeton dans ses
$k^2$ cases.

\subsection*{Problème 3}
Le quadrilatère convexe $ABCD$ est tel que $\widehat{ABC} = \widehat{CDA} = 90^\circ$. Le point $H$ est le pied
de la hauteur issue de $A$ dans le triangle $ABD$. Les points $S$ et $T$ appartiennent respectivement aux côtés
$AB$ et $AD$, de sorte que $H$ soit à l'intérieur du triangle $SCT$ et que
\[
\widehat{CHS} - \widehat{CSB} = 90^\circ, \quad \widehat{THC} - \widehat{DTC} = 90^\circ.
\]
Prouver que la droite $BD$ est tangente au cercle circonscrit au triangle $TSH$.

\bigskip
\begin{center}
  \emph{Second jour --- mercredi 9 juillet 2014 --- Durée : 4 heures 30}
\end{center}

\subsection*{Problème 4}
Les points $P$ et $Q$ appartiennent au côté $BC$ du triangle $ABC$ dont les angles sont aigus, de sorte que
$\widehat{PAB} = \widehat{BCA}$ et $\widehat{CAQ} = \widehat{ABC}$. Les points $M$ et $N$ appartiennent
respectivement aux droites $AP$ et $AQ$, avec $P$ milieu de $[AM]$ et $Q$ milieu de $[AN]$. Prouver que le
point d'intersection des droites $BM$ et $CN$ est sur le cercle circonscrit au triangle $ABC$.

\subsection*{Problème 5}
Pour tout entier $n$ strictement positif, la Banque du Cap émet des pièces de valeur $\frac{1}{n}$. Étant
donné un nombre fini de telles pièces (dont les valeurs ne sont pas nécessairement différentes) de valeur
totale ne dépassant pas $99 + \frac{1}{2}$, prouver qu'il est possible de les répartir en $100$ piles ou
moins, chacune de valeur totale au plus $1$.

\subsection*{Problème 6}
Un ensemble de droites du plan est dit en \emph{position générale} s'il n'y en a pas deux parallèles ni trois
qui passent par un même point. Un ensemble de droites en position générale découpe le plan en régions,
certaines ayant une aire finie. Ces dernières sont appelées les \emph{régions finies de l'ensemble}. Prouver
que pour tout $n$ suffisamment grand, dans tout ensemble de $n$ droites en position générale, il est possible
de colorer en bleu au moins $\sqrt{n}$ de ses droites de sorte qu'il n'y ait aucune région finie de cet
ensemble de bord entièrement bleu.

\emph{Note :} Des résultats pour lesquels $\sqrt{n}$ serait remplacé par $c\sqrt{n}$ seront valorisés selon la
valeur de la constante $c$.
